\section{Introduction}\label{sec:intro}

Let $f$ be a table of $N = 2^n$ field elements indexed by the Boolean hypercube $B_n = \{0,1\}^n$.
Two claims about $f$ recur in succinct arguments: the \emph{hypercube sum}
\[
  S \;=\; \sum_{b \in B_n} f(b),
\]
which is the statement of the sumcheck protocol~\cite{LFKN92}, and the \emph{evaluation} of a multilinear polynomial attached to $f$ at a point $z \in \F^n$, which is the statement of a multilinear polynomial commitment scheme.
Hash-based schemes in the BaseFold paradigm~\cite{ZCF24} commit to a Reed--Solomon encoding of $f$ and prove such claims by interleaving a sumcheck with FRI folding~\cite{BBHR18}, the two sharing their random challenges.

This paper studies the simplest member of this family: the table is committed \emph{as the coefficient vector} of a univariate polynomial, and the scalar claim is reduced by a restriction of one variable per round, with no sumcheck.
Put
\[
  V_f(X) \;=\; \sum_{b=0}^{N-1} f(b)\, X^{b},
  \qquad
  P_f(x_1,\dots,x_n) \;=\; \sum_{b \in B_n} f(b) \prod_{k=1}^{n} x_k^{b_k},
\]
so that $P_f$ is the multilinear polynomial whose \emph{monomial coefficients} are the entries of $f$, and $V_f$ is its Kronecker substitution, $V_f(X) = P_f(X, X^2, \dots, X^{2^{n-1}})$.
Then
\[
  \sum_{b \in B_n} f(b) \;=\; P_f(1,\dots,1) \;=\; V_f(1).
\]
The classical FRI fold $\cfold_\theta(U) = U_e + \theta\, U_o$, where $U(X) = U_e(X^2) + X U_o(X^2)$, substitutes $\theta$ for the first variable of the coefficient-form multilinear polynomial (\cref{sec:restriction}).
For a claim $P_f(z) = v$, the round message of the prover is the affine polynomial
\[
  g_j(T) \;=\; P_{f}(\theta_1,\dots,\theta_{j-1},\, T,\, z_{j+1},\dots,z_n),
\]
the verifier checks $g_j(z_j)$ against the current claim, draws $\theta_j$, and the same $\theta_j$ folds the committed codeword.
For the hypercube sum, $z = (1,\dots,1)$ and $g_j(T) = U_e(1) + T\, U_o(1)$, where $U$ is the current folded polynomial.
After $n$ rounds the folded codeword is the constant $P_f(\theta_1,\dots,\theta_n)$, and the verifier compares it with the reduced claim.
Three states thus move under one challenge sequence: the scalar claim, the coefficient-form polynomial, and the Reed--Solomon word.

The round structure is that of the evaluation protocol of DeepFold~\cite{GLHQTZ25}, which proves coefficient-form evaluations by the same restrictions; DeepFold adds out-of-domain points to reach the list-decoding regime.
The present paper gives the unique-decoding version a complete, self-contained treatment, with the hypercube sum as its distinguished instance, and develops it as the coefficient-form companion of the Boolean-kernel scheme KBFold~\cite{Mkh26}, on the same code and with the same proof machinery.

\subsection{Contributions}\label{sec:intro:contrib}

\paragraph{The restriction dictionary (\cref{sec:restriction}).}
We state the monomial-basis counterpart of the KBFold fold dictionary: the classical fold of $V_f$ is the Kronecker substitution of the restriction of $P_f$ in its first variable, and $n$ folds with $\theta_1,\dots,\theta_n$ give the constant $P_f(\theta_1,\dots,\theta_n)$.
For a word $w$ on a smooth domain $L$, the classical word fold is the line $w_e + \theta\, w_o$ whose generators determine $w$.
This is the property through which correlated agreement enters the soundness proof, and it holds for the classical fold without the reparametrisation of the kernel fold (\KB{Proposition~5.11}).
We record the relation between the two encodings of the same sum: $V_f(1)$ in the monomial basis, and $[X^{N-1}]\,U_f$ in the Boolean-kernel basis.

\paragraph{The scheme and its completeness (\cref{sec:scheme}).}
We define the commitment, the evaluation protocol $\Pieval$ for claims $P_f(z) = v$, and its specialisation $\Pisum$ to hypercube sums, with a stopping parameter, committed levels and the local query checks of KBFold.
The prover sends one affine polynomial per round, that is, one field element once the restriction check $s_j(z_j) = v_{j-1}$ is used to recover the other.
It never forms an equality table, and the commitment requires no change of basis: the table is the coefficient vector.

\paragraph{Soundness (\cref{sec:sound}).}
We prove soundness and evaluation binding in the unique-decoding regime, round-by-round knowledge soundness with a decoding extractor, and the relaxed round-by-round knowledge soundness of Chiesa, Di, Hu and Zheng~\cite{CDHZ25}, which yields post-quantum knowledge soundness of the compiled scheme in the quantum random oracle model.
The proof follows KBFold: the codeword chain and the passing sets carry over unchanged, the folding lemmas are re-proved for the classical fold, and the sumcheck step is replaced by a one-line restriction step.
That step costs $1/|\F|$ per round, against $2/|\F|$ for the degree-$2$ sumcheck of KBFold, and in every round it is dominated by the folding error $|L_j|/|\F|$.
The only external result about codes is again the correlated-agreement theorem of Ben-Sasson, Carmon, Ishai, Kopparty and Saraf~\cite{BCIKS23}.

\paragraph{Formal verification (\cref{app:lean}).}
A Lean~4 formalisation, built on the KBFold Lean library and checked against the same single axiom (correlated agreement), covers the restriction dictionary, completeness, soundness, binding, round-by-round knowledge soundness, the combinatorial part of relaxed round-by-round knowledge soundness, and the transfer to committed levels.

\paragraph{Implementation and experiments (\cref{sec:impl,sec:exp}).}
A Rust implementation reuses the field, NTT, salted Merkle trees and transcript of the KBFold crate.
We compare it on the same machine with four alternatives for the same statements:
\begin{itemize}
  \item the KBFold scheme, which proves the sum as $2^n \tilde f(\tfrac12,\dots,\tfrac12)$;
  \item the coefficient-form scheme with a degree-$2$ sumcheck, as in BaseFold and in the baseline of KBFold;
  \item the quotient opening of $V_f$ at $X = 1$ on a coset domain;
  \item WHIR~\cite{ACFY24}.
\end{itemize}
With the same commitment, verification time and proof size, the opening of a sum is $27$--$41\%$ faster than with KBFold or with the coefficient-form sumcheck scheme, and $34$--$57\%$ faster with post-quantum parameters.

\subsection{Positioning}\label{sec:intro:position}

\Cref{tab:position} lists the protocols that prove hypercube sums or coefficient-form evaluations with univariate encodings, and how each one enters this paper.

\begin{table}[t]
\centering
\small
\begin{tabular}{@{}p{0.2\textwidth}p{0.36\textwidth}p{0.36\textwidth}@{}}
\toprule
Work & Mechanism & Role in this paper \\
\midrule
Sumcheck~\cite{LFKN92} & one variable per round, round polynomial of degree $d$, final evaluation query & the reduction that the restriction step replaces for linear statements (\cref{sec:restriction}) \\
Aurora~\cite{BCRSVW19} & univariate sumcheck over a multiplicative subgroup, for evaluation-form polynomials & different encoding; not compared \\
Gemini~\cite{BCHO22}, PolyFRIM~\cite{ZLGSCLD24} & coefficient-form evaluation by folding at the evaluation point $z$; one commitment per fold & same identity $P_f(z)$ by even--odd splitting; our folds use fresh random challenges and need no extra commitment \\
Zeromorph~\cite{KT24} & the table as univariate coefficients, $\sum_b f(b) X^{b}$; quotients and degree checks with KZG & the same encoding $V_f$ with a pairing-based proof; not hash-based \\
BaseFold~\cite{ZCF24} & FRI folds whose last constant is a random evaluation; degree-$2$ sumcheck on $f\cdot\eq_z$ & source of the final-value identity; the coefficient-form sumcheck baseline of \cref{sec:exp} \\
DeepFold~\cite{GLHQTZ25} & coefficient-form evaluation by restrictions sharing the FRI challenges; out-of-domain points for list decoding & the round structure of $\Pieval$; we treat its unique-decoding version, with full proofs \\
FriEval~\cite{Zha25} & FRI as a protocol evaluating a committed multilinear at challenge points & the same reading of FRI, in a functional framework \\
Fold-DCS~\cite{LMN25} & sumcheck in $O(\log n)$ rounds by divide and conquer & different technique; not compared \\
WHIR~\cite{ACFY24}, Hab\"ock~\cite{Hab24} & constrained Reed--Solomon codes; list decoding via subcode correlated agreement & external baseline; route to the list-decoding version (\cref{sec:concl}) \\
KBFold~\cite{Mkh26} & table encoded in the Boolean-kernel basis; kernel fold evaluates $\tilde f$ & companion scheme: same code, soundness machinery and Lean library \\
\bottomrule
\end{tabular}
\caption{Protocols for hypercube sums and coefficient-form evaluations with univariate encodings.}
\label{tab:position}
\end{table}

The coefficient-form polynomial $P_f$ is not the multilinear extension $\tilde f$ of the table: $P_f(b) = \sum_{c \preceq b} f(c)$ on $B_n$.
A system whose statements concern $\tilde f$, such as a sumcheck-based argument that ends with a claim $\tilde f(r) = v$, can use $\Pieval$ after the M\"obius transform of the table, which is what BaseFold and DeepFold do, or use KBFold, which avoids the transform.
The hypercube sum is the one statement on which the two polynomials agree without any transform, since $\sum_b f(b) = P_f(1,\dots,1)$ and, in odd characteristic, $\sum_b f(b) = 2^n \tilde f(\tfrac12,\dots,\tfrac12)$.

\subsection{Related work}\label{sec:intro:related}

\paragraph{Proximity tests and their soundness.}
FRI~\cite{BBHR18} tests proximity to Reed--Solomon codes by repeated even--odd folding, and its soundness rests on proximity gaps for Reed--Solomon codes~\cite{BCIKS23}.
DEEP-FRI~\cite{BGKS20} samples outside the evaluation domain to bind a word to one codeword in the list-decoding regime, the technique DeepFold uses for multilinear evaluations.
WHIR~\cite{ACFY24} tests proximity to constrained Reed--Solomon codes with very fast verification.
Hab\"ock~\cite{Hab24} proves BaseFold sound up to the Johnson bound through correlated agreement for linear subcodes.
Zhang~\cite{Zha25} recasts FRI as a functional proof system evaluating committed multilinear polynomials and builds the security of BaseFold on the theorems of FRI.
Our analysis stays in the unique-decoding regime and uses only the correlated-agreement theorem.

\paragraph{Multilinear commitments from folding.}
BaseFold~\cite{ZCF24} introduced the interleaving of sumcheck rounds with folding rounds under shared challenges and the observation that the last FRI constant is a random evaluation of the committed coefficient-form polynomial.
DeepFold~\cite{GLHQTZ25} replaces the sumcheck by restrictions for coefficient-form polynomials and reaches the list-decoding regime; it also batches polynomials of different sizes and adds zero knowledge.
KBFold~\cite{Mkh26} encodes the table in the Boolean-kernel basis so that the folds evaluate the multilinear extension of the table itself.

\paragraph{Multilinear-to-univariate reductions.}
Gemini~\cite{BCHO22} reduces a coefficient-form claim $P(z) = v$ to univariate claims by folding at the coordinates of $z$ and committing to every folded polynomial; PolyFRIM~\cite{ZLGSCLD24} instantiates this approach with FRI.
Zeromorph~\cite{KT24} maps the table to univariate coefficients, which is the encoding $V_f$ of this paper, and proves evaluations of $\tilde f$ with quotients and degree checks.
Mohamed~\cite{Moh25} classifies these adaptors by how they map multilinear to univariate polynomials and gives a values-to-values reduction.
These works treat the univariate commitment as a black box; here the univariate polynomial is committed through a Reed--Solomon code whose own folding carries the reduction.

\paragraph{Sumcheck variants.}
The sumcheck protocol~\cite{LFKN92} reduces a hypercube sum to one evaluation, with one round polynomial of degree $d$ per variable.
Aurora~\cite{BCRSVW19} proves sums over multiplicative subgroups for univariate polynomials in one round.
Fold-DCS~\cite{LMN25} reduces the number of rounds to $O(\log n)$ by a divide-and-conquer recursion.
For the linear statement $\sum_b f(b) = S$ about a committed coefficient vector, the restriction step of this paper plays the role of the sumcheck, with affine round polynomials and no final evaluation query.

\subsection{Organisation}\label{sec:intro:org}

\Cref{sec:prelim,sec:proofs} recall, without proofs, the notation and the results of KBFold that we use: multilinear polynomials, smooth domains, Reed--Solomon codes, correlated agreement, interactive oracle proofs, round-by-round soundness and the theorem of~\cite{CDHZ25}.
\Cref{sec:restriction} proves the restriction dictionary and the line form of the classical fold.
\Cref{sec:scheme} defines the scheme and proves completeness, and \cref{sec:sound} proves soundness, evaluation binding, round-by-round knowledge soundness and post-quantum knowledge soundness.
\Cref{sec:impl,sec:exp} describe the implementation and the measurements, and \cref{sec:concl} concludes.
\Cref{app:lean} maps the paper to the Lean formalisation.

Results of KBFold are cited by their numbers in~\cite{Mkh26}, as in \KB{Theorem~2.23}.
Every result specific to this paper is proved here.
